% =====================================================================
%  THÈME 4 — Logarithmes, pH & incertitude — CORRIGÉ — Niveau FACILE
% =====================================================================
\documentclass[11pt,a4paper]{article}
\usepackage[utf8]{inputenc}
\usepackage[T1]{fontenc}
\usepackage{lmodern}
\usepackage[french]{babel}
\usepackage{amsmath,amssymb,mathtools}
\usepackage{icomma}
\usepackage{siunitx}
\sisetup{output-decimal-marker={,},per-mode=symbol}
\usepackage[version=4]{mhchem}
\usepackage{xcolor}
\usepackage[most]{tcolorbox}
\usepackage{enumitem}
\usepackage{fancyhdr}
\usepackage[a4paper,margin=2.2cm,headheight=26pt,headsep=8pt]{geometry}

\definecolor{primary}{RGB}{150,98,162}
\definecolor{primarydark}{RGB}{92,52,110}
\newcommand{\cs}{chiffre significatif}
\newcommand{\CS}{chiffres significatifs}

\pagestyle{fancy}\fancyhf{}
\fancyhead[L]{\small\textcolor{primarydark}{\textbf{Fiche 1} — Chiffres significatifs}}
\fancyhead[R]{\small\textcolor{primarydark}{\textsc{Corrigé · Facile · Thème 4}}}
\fancyfoot[C]{\small\thepage}
\renewcommand{\headrulewidth}{0.8pt}
\renewcommand{\headrule}{\hbox to\headwidth{\color{primary}\leaders\hrule height \headrulewidth\hfill}}

\newtcolorbox[auto counter]{cor}[1][]{enhanced,breakable,boxrule=0.8pt,arc=2pt,
  colback=primarydark!5,colframe=primarydark,left=3mm,right=3mm,top=2mm,bottom=2mm,
  fonttitle=\bfseries,coltitle=white,
  title={Correction — Exercice~\thetcbcounter\ifx\relax#1\relax\relax\else\ ---\ \textmd{\itshape #1}\fi}}

\setlength{\parindent}{0pt}\setlength{\parskip}{3pt}

\begin{document}

\begin{center}
{\Large\bfseries\color{primarydark} Corrigé — Niveau Facile}\\[1pt]
{\color{primary} Thème 4 : pH, justesse, fidélité, incertitude}
\end{center}
\vspace{2mm}

\begin{cor}[Combien de décimales pour le pH ?]
Règle : nombre de décimales du pH $=$ nombre de \CS\ de $[\ce{H3O+}]$.
\begin{itemize}[leftmargin=1.6em,itemsep=1pt]
  \item a) $2$~CS $\to \mathbf{2}$ décimales. \quad b) $3$~CS $\to \mathbf{3}$ décimales.
  \item c) $2$~CS $\to \mathbf{2}$ décimales. \quad d) $2$~CS $\to \mathbf{2}$ décimales.
\end{itemize}
\end{cor}

\begin{cor}[Du pH à la concentration (valeurs simples)]
Le nombre de décimales du pH donne le nombre de \CS\ de la concentration.
\begin{itemize}[leftmargin=1.6em,itemsep=1pt]
  \item a) $\mathrm{pH}=\num{2,00} \to [\ce{H3O+}]=\num{1,0e-2}\,\si{\mol\per\litre}$ ($2$~CS).
  \item b) $\mathrm{pH}=\num{5,00} \to \num{1,0e-5}\,\si{\mol\per\litre}$ ($2$~CS).
  \item c) $\mathrm{pH}=\num{7,000} \to \num{1,00e-7}\,\si{\mol\per\litre}$ ($3$~CS).
  \item d) $\mathrm{pH}=\num{11,0} \to \num{1e-11}\,\si{\mol\per\litre}$ ($1$~CS).
\end{itemize}
\end{cor}

\begin{cor}[Incertitude absolue sous-entendue]
$\Delta x$ porte sur le rang du dernier \cs\ écrit.
\begin{itemize}[leftmargin=1.6em,itemsep=1pt]
  \item a) $\Delta m=\qty{0,01}{\gram}$. \quad b) $\Delta V=\qty{0,01}{\milli\litre}$.
  \item c) $\Delta m=\qty{0,1}{\gram}$. \quad d) $\Delta L=\qty{0,001}{\metre}$.
\end{itemize}
\end{cor}

\begin{cor}[Justesse ou fidélité ?]
\begin{itemize}[leftmargin=1.6em,itemsep=1pt]
  \item a) groupés mais décentrés : \textbf{fidèles mais non justes} (reproductible, mais biaisé).
  \item b) dispersés autour du centre : \textbf{justes (en moyenne) mais non fidèles}.
  \item c) groupés au centre : \textbf{justes et fidèles}.
\end{itemize}
\end{cor}

\begin{cor}[Encadrer une mesure]
\begin{itemize}[leftmargin=1.6em,itemsep=1pt]
  \item a) $m\in[\,\qty{13,01}{}\,;\ \qty{13,03}{\gram}\,]$.
  \item b) $V\in[\,\qty{20,95}{}\,;\ \qty{20,97}{\milli\litre}\,]$.
\end{itemize}
\end{cor}

\end{document}
